\documentclass[xcolor=dvipsnames,notheorems,aspectratio=169]{beamer}
%\documentclass[xcolor=dvipsnames,notheorems,aspectratio=169]{beamer}
\input{preamble}
\input{quote}
\input{style}
%\input{risc-v}

\usepackage{listings}
\usepackage[most]{tcolorbox}
\usepackage[table]{xcolor}
\definecolor{lavender}{rgb}{0.85, 0.88, 0.9}
\definecolor{pastelred}{rgb}{1.0, 0.8, 0.8}
\definecolor{pastelgreen}{rgb}{0.8, 1.0, 0.8}
\definecolor{pastelorange}{rgb}{1.0, 0.9, 0.8}
\definecolor{pastelblue}{rgb}{0.8, 0.87, 1.0}
\definecolor{pastelslate}{rgb}{0.85, 0.88, 0.9}

\definecolor{bluetext}{RGB}{38, 71, 176}
\definecolor{orangetext}{RGB}{222, 114, 49}
\definecolor{grayrow}{gray}{0.95}

\usepackage{algorithm}
\usepackage{algpseudocode}
\usepackage{booktabs}

\newif\ifshowsolutions
\showsolutionstrue
%\showsolutionsfalse

\newtcolorbox[auto counter, number within=section]{bluebox}[1]{
    colback=blue!10, colframe=blue!80!black,
    coltitle=white, fonttitle=\bfseries,
    title={#1}
}

\newtcolorbox[auto counter, number within=section]{redbox}[2][]{
    colback=red!10, colframe=red!80!black,
    coltitle=white, fonttitle=\bfseries,
    title={#2}
}

\newtcolorbox[auto counter, number within=section]{greenbox}[2][]{
    colback=green!10, colframe=green!80!black,
    coltitle=white, fonttitle=\bfseries,
    title={#2}
}

\newtcolorbox[auto counter, number within=section]{solutionbox}[2][]{
    colback=gray!10, colframe=gray,
    coltitle=white, fonttitle=\bfseries,
    title={#1}
}

\newtcolorbox[auto counter, number within=section]{cyanbox}[2][]{
    colback=cyan!8,
    colframe=cyan!60!black,
    coltitle=white,
    fonttitle=\bfseries,
    title={#2},
    arc=8pt
}

\newtcolorbox[auto counter, number within=section]{orangebox}[2][]{
    colback=orange!10,
    colframe=orange!70!black,
    coltitle=white,
    fonttitle=\bfseries,
    title={#2},
    arc=8pt
}

\newtcolorbox[auto counter, number within=section]{lemmabox}[1][]{
    colback=purple!10,
    colframe=purple!70!black,
    coltitle=white,
    fonttitle=\bfseries,
    title={#1},
    arc=8pt
}





% Pauses in math environment
\makeatletter
\let\save@measuring@true\measuring@true
\def\measuring@true{%
  \save@measuring@true
  \def\beamer@sortzero##1{\beamer@ifnextcharospec{\beamer@sortzeroread{##1}}{}}%
  \def\beamer@sortzeroread##1<##2>{}%
  \def\beamer@finalnospec{}%
}
\makeatother

\usetikzlibrary{matrix}
\usetikzlibrary{trees}

\providecommand{\tape}[3]{
  \foreach \x [count=\y from 0] in {1,...,#2}{
    \ifnum\y<1
        \if#3\empty
            \node (#1-\x) [minimum height=1.5em,minimum width=1.5em,outer sep=0pt,draw,rectangle,node distance=0pt] {};
        \else
            \node (#1-\x) [minimum height=1.5em,minimum width=1.5em,outer sep=0pt,draw,rectangle,node distance=0pt] at #3 {};
        \fi
    \else
        \node (#1-\x) [minimum height=1.5em,minimum width=1.5em,outer sep=0pt,draw,rectangle,node distance=0pt, right=of #1-\y] {};
    \fi
  }
}

\lstdefinelanguage[RISC-V]{Assembler}
{
  alsoletter={.}, % allow dots in keywords
  alsodigit={0x}, % hex numbers are numbers too!
  morekeywords=[1]{ % instructions
    lb, lh, lw, li, lbu, lhu,
    sb, sh, sw,
    sll, slli, srl, srli, sra, srai,
    add, addi, sub, lui, auipc,
    xor, xori, or, ori, and, andi,
    slt, slti, sltu, sltiu,
    beq, bne, blt, bge, bltu, bgeu,
    j, jr, jal, jalr, ret,
    scall, break, nop
  },
  morekeywords=[2]{ % sections of our code and other directives
    .align, .ascii, .asciiz, .byte, .data, .double, .extern,
    .float, .globl, .half, .kdata, .ktext, .set, .space, .text, .word
  },
  morekeywords=[3]{ % registers
    zero, ra, sp, gp, tp, s0, fp,
    t0, t1, t2, t3, t4, t5, t6,
    s1, s2, s3, s4, s5, s6, s7, s8, s9, s10, s11,
    a0, a1, a2, a3, a4, a5, a6, a7,
    ft0, ft1, ft2, ft3, ft4, ft5, ft6, ft7,
    fs0, fs1, fs2, fs3, fs4, fs5, fs6, fs7, fs8, fs9, fs10, fs11,
    fa0, fa1, fa2, fa3, fa4, fa5, fa6, fa7
  },
  morecomment=[l]{;},   % mark ; as line comment start
  morecomment=[l]{\#},  % as well as # (even though it is unconventional)
  morestring=[b]",      % mark " as string start/end
  morestring=[b]'       % also mark ' as string start/end
}

% usage example:

% define some basic colors
\definecolor{mauve}{rgb}{0.58,0,0.82}

\lstset{
  % listings sonderzeichen (for german weirdness)
  literate={ö}{{\"o}}1
           {ä}{{\"a}}1
           {ü}{{\"u}}1,
  basicstyle=\normalsize\ttfamily,                    % very small code
  breaklines=true,                              % break long lines
  commentstyle=\itshape\color{green!50!black},  % comments are green
  keywordstyle=[1]\color{blue!80!black},        % instructions are blue
  keywordstyle=[2]\color{orange!80!black},      % sections/other directives are orange
  keywordstyle=[3]\color{red!50!black},         % registers are red
  stringstyle=\color{mauve},                    % strings are from the telekom
  identifierstyle=\color{teal},                 % user declared addresses are teal
  frame=l,                                      % black line on the left side of code
  language=[RISC-V]Assembler,                   % all code is RISC-V
  tabsize=4,                                    % indent tabs with 4 spaces
  showstringspaces=false                        % do not replace spaces with weird underlines
}

\renewcommand{\S}{\mathcal{S}}

%\usepackage{enumitem}
% Set spacing for all itemize environments
%\setlist[itemize]{before=\vskip 0.5cm, after=\vskip 0.5cm}

\makeatletter
%\g@addto@macro\normalsize{%
%    \setlength\belowdisplayskip{-0pt}
%}

\AtBeginSection[]
{
    \begin{frame}
        \usebeamerfont{title}\insertsectionhead\par%
    \end{frame}
}



\newcommand{\inserteqstrut}[1]{%
\rlap{\ensuremath{#1}}%
\phantom{\hspace{10cm}}}
    
\providecommand{\ito}{It\^o}

\title{ENG17: Circuits I \\ Circuit Analysis Techniques: Thevenin and Norton Equivalence + Maximum Power Transfer}
\author{Anthony Thomas}

\begin{document}

\begin{frame}
    \maketitle
\end{frame}

\begin{frame}{Abstraction}
\begin{itemize}
    \item Modern engineering problems are extraordinarily complex:
\end{itemize}
\begin{center}
    \includegraphics[width=0.8\linewidth]{figures/phone-block-diagram.png}
\end{center}
\end{frame}

\begin{frame}{Abstraction}
\begin{itemize}
    \item Even individual pieces are highly sophisticated\footnote{ADC block diagram for TI-MSP432P microcontroller line.}:
\end{itemize}
\begin{center}
    \includegraphics[width=0.6\linewidth]{figures/precision-adc-block-diagram.png}
\end{center}
\end{frame}

\begin{frame}{Abstraction}
\begin{cyanbox}{Modularity}
\begin{itemize}
    \item Our ability to design modern engineered systems is completely predicated on being able to break them down into simpler pieces that can be developed (more or less) independently.
    \item Each component provides an \emph{interface} that allows us to understand that component's functionality and behavior without knowing \emph{exactly} how it is implemented.
    \item I don't want to have to understand the low-level details of RF design to work on processor microarchitecture.
\end{itemize}
\end{cyanbox}
\end{frame}

\begin{frame}{Thevenin's Theorem}
\begin{itemize}
    \item Thevenin's and Norton's Theorems are fundamental results in this regard:
\end{itemize}
\begin{orangebox}{Thevenin's Theorem (Thevenin Equivalence)}
    \begin{itemize}
        \item A linear circuit can be represented at its output terminals by an equivalent circuit consisting of a series combination of a voltage source $v_{th}$ and a resistor $R_{th}$.
    \end{itemize}
\end{orangebox}
\begin{greenbox}{Intuition}
\begin{itemize}
    \item Thevenin Equivalence allows us to treat potentially very complex portions of a design as a simple equivalent circuit that ``appears the same'' to the rest of the circuit.
    \item By ``appears the same'' we mean: delivers the same current and has the same voltage across its terminals.
\end{itemize}
\end{greenbox}
\end{frame}

\begin{frame}[t]{Finding $v_{th}$}
\begin{itemize}
    \item Finding the Thevenin equivalent circuit amounts to finding: $v_{th}$ and $R_{th}$.
    \item \textbf{Finding}: $v_{th}$
    \begin{itemize}
        \item For equivalence to hold, no matter what ``load'' ($R_{L}$) we connect, both the original circuit and the Thevenin equivalent must supply the same current $i_{L}$.
        \item How can we use this to calculate $v_{th}$?
    \end{itemize}
\end{itemize}
\end{frame}

\begin{frame}{Finding $R_{th}$}
\begin{itemize}
    \item Three common methods to find $R_{th}$:
    \begin{itemize}
        \item[(1)] Short-Circuit Method
        \item[(2)] Equivalent-Resistance Method
        \item[(3)] External-Source Method
    \end{itemize}
    \item Differ in assumptions about nature of voltage/current sources. 
\end{itemize}
\end{frame}

\begin{frame}{Norton Equivalence}
\begin{itemize}
    \item By applying source-transformation to the Thevenin equivalent, we obtain the following \emph{corollary}:\footnote{A corollary is a result that follows as a more-or-less straightforward consequence of a theorem.}
\end{itemize}
\begin{orangebox}{Norton Equivalence}
\begin{itemize}
    \item A linear circuit can be represented at its output terminals by an equivalent circuit consisting of a parallel combination of a current source $i_{N}$ and a resistor $R_{N}$.
\end{itemize}
\end{orangebox}
\begin{redbox}{Test Yourself}
\begin{itemize}
    \item Given a circuit in Thevenin equivalent form, derive the Norton equivalent.
\end{itemize}
\end{redbox}
\end{frame}

\begin{frame}{Maximum Power Transfer}
\begin{itemize}
    \item A famous electrical engineer:
\end{itemize}
\begin{center}
    \includegraphics[width=0.75\linewidth]{figures/unlimited-power.png}
\end{center}
\end{frame}

\begin{frame}[t]{Maximum Power Transfer}
\begin{itemize}
    \item In real life, power is limited, but we may want to maximize the power transferred from a source to some ``load'' (e.g. resistor).
    \item \textbf{Question}: for what load resistance is the power delivered to $R_{L}$ maximized?
    \pause
    \item Perhaps surprisingly:
\end{itemize}
\begin{cyanbox}{Maximum Power Transfer}
\begin{itemize}
    \item For an active (power-generating) circuit modeled by a Thevenin equivalent $v_{s}$ and $R_{s}$, power transfer to a passive load $R_{L}$ is maximized when:
    \[
        R_{L} = R_{s},
    \]
    and the power delivered is:
    \[
        p_{L}^{*} = \frac{v_{s}^{2}}{4R_{L}}.
    \]
\end{itemize}
\end{cyanbox}
\end{frame}

\begin{frame}[t]{Maximum Power Transfer}
\begin{itemize}
    \item Where did this come from?
\end{itemize}
\end{frame}

\end{document}
